\chapter{Introducción}
\label{cap:introduccion}

En una agencia de turismo es importante tener bien organizada la
información de los viajes que ofrece. Cada paquete tiene un precio, un
itinerario, fechas de salida y una cantidad determinada de lugares
disponibles. Cuando un cliente solicita información o desea hacer una
reserva, el personal necesita revisar estos datos para darle una respuesta
correcta y confirmar que todavía existen cupos.

La ciudad de Potosí cuenta con atractivos culturales y naturales que
generan interés turístico. Uno de los más representativos es la ciudad
histórica, inscrita en la Lista del Patrimonio Mundial en 1987
\parencite{unescoPotosi}. Además, el Instituto Nacional de Estadística
publica información sobre el flujo de visitantes internacionales en Bolivia
\parencite{ineVisitantes}. En este contexto, resulta pertinente estudiar
cómo las agencias de turismo de la ciudad organizan sus paquetes y atienden
las solicitudes de reserva.

El problema central que se plantea es la ineficiencia en el proceso de
venta y reserva de paquetes turísticos. Esta se manifiesta, entre otros
aspectos, en la verificación no uniforme de cupos, en la actualización
desfasada de la información entre canales de atención y en la falta de
criterios claros para modificar, cancelar y liberar reservas. El problema
no radica en que la agencia utilice cuadernos, hojas de cálculo u otros
medios de registro, sino en que, al verificar, modificar o confirmar una
reserva, todo el personal no trabaje con la misma información. Cuando esto
ocurre, pueden aparecer diferencias entre lo que se ofrece al cliente y los
lugares que realmente quedan disponibles.

Por ello, se propone desarrollar un sistema web que permita administrar
paquetes, fechas de salida y reservas desde una misma plataforma, con
funciones de apoyo como el registro de clientes, las cotizaciones, los
pagos, los proveedores y los reportes. La propuesta se presenta como un
proyecto académico, y sus causas se comprobarán mediante entrevistas y
observación en la agencia que se seleccione.
